%% ============================================================
%% 分片 30：原文第 10 节「The Bogomolny Equations And The Space Of Hecke Modifications」（原文 glshort.tex 起于第 8120 行）
%% 由 tools/harvmac2tex.py 从 arXiv 源码（hep-th/0604151v3）自动生成。
%% 只译英文散文；公式、编号、\label、\cite 一律不改。
%% ============================================================

\subsection{扩展的 Bogomolny 方程}

\label{sec:extbog}

迄今为止，在研究 't Hooft 算子对膜的作用时，我们只考虑过不依赖于时间且时间反演不变的超对称解。在这一限制下，BPS 方程 \eqref{eq:oops}\ 约化为普通的 Bogomolny 方程，而 't Hooft 算子约化为通常的几何 Hecke 算子。

不依赖时间的假设其实一般来说是可以论证的，但要求时间反演不变性会对允许的边界条件施加严重限制。只有当我们挑选膜的方式（正如我们在第~\ref{sec:opprod}~节讨论 't Hooft 算子的乘积时所做的）能保证相关解满足 $\varphi=0$ 时，这个要求才是恰当的。对几何 Langlands 对偶而言最重要的那些 $A$-膜（Hitchin 纤维化的一般纤维，以及第~\ref{sec:abranes}~节所描述的 c.c. 膜），其支撑并不限于 $\varphi=0$；因此，为了研究 't Hooft 算子对它们的作用，我们需要放弃时间反演不变性这一条件。这正是下面讨论的目的。远为详尽的内容将另文给出 \cite{ref26}。

我们再次考虑四维流形 $M=\RR\times W$ 上的 BPS 方程，并把 Higgs 场 $\phi$ 写成 $\phi_0\, dx^0 +\pi^*\tilde\phi$，其中 $\pi$ 是到 $W$ 的投影，$\tilde\phi$ 是 $W$ 上的一形式 Higgs 场。同样，我们把规范场写成 $A=A_0\,dx^0+\tilde A$，其中 $\tilde A$ 是曲率为 $\tilde F$ 的三维联络。$t=1$ 时，不依赖于时间的 BPS 方程为
% 式 (10.35)
\setcounter{equation}{34}
\begin{equation}\label{eq:BPSstaticgen}
\begin{aligned}\tilde F-\tilde\phi\wedge\tilde\phi&=
\star\left(D\phi_0-[A_0,\tilde\phi]\right),\\ \star
D\tilde\phi&=[\phi_0,\tilde\phi]+D A_0,\\
D^*\tilde\phi+[A_0,\phi_0]&=0.\end{aligned}
\end{equation}
 这里外微分 $D$、Hodge 星算子 $\star$ 以及算子 $D^*=\star D\star$ 都按三维意义理解。

这些方程借助适当的消灭定理可以大大简化。事实上，用与时间反演不变情形相同的论证和同类型的假设（回忆式~\eqref{eq:wungo}），可以证明 $\phi_1=A_0=0$。

于是我们令 $\phi_1=A_0=0$，在 $C$ 上取局部复坐标 $z=x^2+ix^3$，记 $y=x^1$，并取规范 $\tilde A_y=0$。我们从此还把记号中的波浪号去掉。并且，正如第~\ref{sec:theck}~节的做法，我们取 $W$ 上的度规为 $ds^2=h |dz|^2+ dy^2$，其中 $h$ 是某个函数。于是扩展的 Bogomolny 方程 \eqref{eq:BPSstaticgen}\ 约化为
% 式 (10.36)
\setcounter{equation}{35}
\begin{equation}\label{eq:BPSextsimple}
\begin{aligned}D_\bz\phi_z&=0,\\
\partial_y A_\bz&=-i D_\bz \phi_0,\\
\partial_y\phi_z&=-i[\phi_z,\phi_0],\\
F_{z\bz}-[\phi_z,\phi_\bz]&={ih\over 2}\partial_y\phi_0.\end{aligned}
\end{equation}

\eqref{eq:BPSextsimple} 的第一式表明：对任意 $y\in I$，限制在 $C_y=\{y\}\times C$ 上的 $\varphi=\phi_z\,dz$ 是 $\End(E)\otimes K_C$ 的全纯截面。于是，若把由 $\bar\partial$ 算子 $\bar D=d\bar z(\partial_{\bar z}+A_{\bar z})$ 在 $C_y$ 上定义的全纯丛记作 $E$，则对任意 $y$，$(E,\varphi)$ 都是一个 Higgs 丛。接下来两式的解释，根据我们在第~\ref{sec:theck}~节处理普通 Bogomolny 方程的经验，应当不难预料：它们说的是，沿 $y$ 方向的演化只以规范变换的方式改变 Higgs 对 $(E,\varphi)$。因此，在遇到与 't Hooft 算子相联系的奇点之前，扩展 Bogomolny 方程的一个解所决定的 $\MH$ 中的点与 $y$ 无关。

现在设有一个权为 $^L\negthinspace w$ 的 't Hooft 算子，位于（比方说）$y=y_0$、$z=z_0$。当 $y$ 增大越过 $y=y_0$ 时，丛 $E$ 经历一次普通的 Hecke 修正，正如第~\ref{sec:theck}~节所描述的那样。事实上，$A_{\bar z}$ 的演化方程与第~\ref{sec:theck}~节中的相同，由 't Hooft 算子引起的奇点也相同。对 $(E,\varphi)$ 同样发生我们将称为 Hecke 修正的变化。然而 $\varphi$ 的作用与 $E$ 的作用大不相同。

全纯丛可以在单个点上被修正而保持全纯性，Hecke 算子或 't Hooft 算子所做的正是这件事。这种可能的修正构成一个模空间，即大家熟悉的、参数化 Bogomolny 方程奇异解的空间 $\bar {\EUY}({}^L\negthinspace w;z_0)$（它同构于一个 Schubert 簇）。但一个全纯场——在此情形是 ${\rm ad}(E)\otimes K_C$ 的全纯截面 $\varphi$——不能在一点被修正而保持全纯性。因此，如果 $\varphi$ 在 $z=0$ 之外给定，那么越过 $z=0$ 它至多只有唯一一个全纯延拓。

所以，把 $\varphi$ 包括进来，并不会使 't Hooft 算子引起的 Hecke 修正依赖于额外的参数。恰恰相反，一般地，包括 $\varphi$ 会消去一些参数。原因在于：即使 $\varphi$ 在 $y<y_0$ 时于 $z=z_0$ 处全纯，在 $y>y_0$ 时它一般也会在 $z=z_0$ 处出现极点。为看清这一点，我们取 $G=U(N)$，考虑权为 $^L\negthinspace w=(m_1,\ldots,m_N)$ 的 't Hooft 算子。越过 $y=y_0$ 时丛 $E$ 的全纯型发生跳变；我们把「之前」与「之后」的丛分别记为 $E_-$ 和 $E_+$。相对于某个局部分解 $E_-\cong \oplus_{i=1}^N \cL_i,$ 丛 $E_+$ 有局部分解 $\oplus_{i=1}^N\cL_i\otimes \CO(z_0)^{m_i}$。也就是说，若 $\cL_i$ 局部由全纯截面 $s_i$ 生成（$i=1,\ldots,N$），则 $E_+$ 由全纯截面 % 式 (10.37)
\setcounter{equation}{36}
\begin{equation}\label{eq:zilfog}
t_i=(z-z_0)^{-m_i} s_i.
\end{equation}
 这里我们对 $z\neq 0$ 利用了 $E_+$ 与 $E_-$ 之间的一个自然同构。

设 $\varphi$ 在 $y<y_0$ 时于 $z_0$ 处全纯。乘积 $\varphi \,s_i$ 是 $E_- \otimes K_C$ 的全纯截面（因为可以在 $z=z_0$ 附近把它平凡化，$K_C$ 在此不起本质作用），故可按 $s_i$ 展开：
% 式 (10.38)
\setcounter{equation}{37}
\begin{equation}\label{eq:bongo}
\varphi\cdot s_i=\sum_j f_i{}^j\,s_j.
\end{equation}
  $\varphi$ 在 $y<y_0$ 时于 $z=z_0$ 处全纯，这一事实说明 $f_i{}^j$（它们是 $K_C$ 的截面）在 $z=z_0$ 处全纯。它们不再受其他约束。

同样，对 $y>y_0$，我们可以把 $\varphi\,t_i$ 展开为 $t_j$ 的线性组合。给定 $s_i$ 与 $t_i$ 之间的关系 \eqref{eq:zilfog}\ 后，我们可以立即写出展开式的形式： % 式 (10.39)
\setcounter{equation}{38}
\begin{equation}\label{eq:bzongo}
\varphi\cdot t_i=\sum_j
(z-z_0)^{m_j-m_i}f_i{}^j\,t_j.
\end{equation}
 要使 $\varphi$ 在 $y>y_0$ 时全纯，函数 $(z-z_0)^{m_j-m_i}f_i{}^j$ 必须在 $z=z_0$ 处全部正则。若把 $m_i$ 按通常方式排序为 $m_1\geq m_2\geq\dots\geq m_N$，则对 $f_i{}^j$ 在 $z=z_0$ 附近的 Taylor 级数系数所要施加的条件数目为 $\sum_{i<j}(m_i-m_j)$。

这与式~\eqref{eq:ydro} 所确定的丛的 Hecke 修正空间的维数相同。这强烈提示：与丛的 Hecke 修正情形形成鲜明对照，Higgs 丛 $(E,\varphi)$ 的 Hecke 修正一般构成有限集。事实上，只要 $\varphi(z_0)$ 正则且半单（对 $G=U(N)$ 而言即它可对角化且本征值互异），情形就是如此。对结果最系统的描述是：$\varphi$ 在参数化 $E$ 于 $z_0$ 处 $^L\negthinspace w$ 型 Hecke 修正的簇 $\bar {\EUY}({}^L\negthinspace w,z_0)$ 上确定一个向量场；该向量场的零点就是 $(E,\varphi)$ 的可能的 Hecke 修正。（这个向量场并不完全由典范方式确定，但它的零点是。）

把详细解释留待 \cite{ref26}，我们简要描述 $G=U(2)$、$^L\negthinspace w=(1,0)$ 的例子。在这个例子中，正如我们在第~\ref{sec:spaceheck}~节所学到的，若取丛 $E_-$ 为平凡丛，则 $E_+$ 可以这样刻画：$E_+$ 在 $z=0$ 处全纯的截面具有形式 % 式 (10.40)
\setcounter{equation}{39}
\begin{equation}\label{eq:hxonxno}
s_+=(z-z_0)^{-1}
u+(h_1(z),h_2(z)),
\end{equation}
  其中 $u$ 是 $\Bbb{C}^2$ 的某个非零元素，$h_1$ 与 $h_2$ 在 $z=z_0$ 附近全纯。在模去 $\Bbb{C}^*$ 的伸缩后，$u$ 参数化了 $E$ 的可能的 Hecke 修正所构成的空间，在本例中它就是 $\Bbb{CP}^1$。$\varphi$ 在 $y>y_0$ 时于 $z=z_0$ 处全纯，当且仅当乘积 $\varphi\cdot s_+$ 具有同样的形式，换言之，当且仅当 $\varphi u$ 是 $u$ 的倍数。在正则半单情形，这意味着 $u$ 必须是 $\varphi$ 的两个特征向量之一。

回到扩展的 Bogomolny 方程 \eqref{eq:BPSextsimple}，由关系 $\partial_y\phi_z=-i[\phi_z,\phi_0]$ 立即得到：若 $\phi_z$ 在某个 $y$ 处为零，则它处处为零。我们在第~\ref{sec:opprod}~节分析 't Hooft 算子的乘积时假设了这一点。同一方程的又一个推论是：$\varphi$ 的特征多项式与 $y$ 无关，即使在越过一个 't Hooft 算子时也是如此。因此，Higgs 丛的 Hecke 变换把 Hitchin 纤维化的每根纤维映入它自身。

现在考虑一根缠在 Hitchin 纤维化某根纤维上、并带有平坦 Chan-Paton 丛的膜。按照 $S$-对偶，这类膜必是磁特征膜。由上一段所述，这种膜的支撑在 Hecke 变换下不变。这一观察是向着不依赖 $S$-对偶而直接证明这类膜是磁特征膜迈出的一步，细节将另文给出 \cite{ref26}。

\par\medskip{\it Higgs 丛的 Hecke 对应}

正如我们在讨论普通 Bogomolny 方程时那样，我们可以考虑给定 Higgs 丛 $(E,\varphi)$ 的可能的 Hecke 修正——这也是我们迄今的观点——也可以更对称地描述初态与末态，即描述 Higgs 丛空间上的一个 Hecke 对应 ${\EUQ}$。我们这样来定义 ${\EUQ}$：在 $W=I\times C$ 上求解扩展 Bogomolny 方程，在两端取适当的边界条件，并在指定 't Hooft 算子的存在下求解。我们在 $W$ 的两端同时指定 $A$ 与 $\varphi$（但不指定 $\phi_0$），并且只模去在每一端均为平凡的规范变换。

设 $\EUW_0$ 为 $C$ 上满足 $\bar D\varphi=0$ 的规范场与 Higgs 场构成的全体对 $(A,\varphi)$ 的空间。那么 $\EUW_0$ 就是全体 Higgs 丛的「栈（stack）」${\rm Higgs}_G$ 的规范场论版本，正如全体联络的空间是全体 $G$-丛的栈 ${\rm Bun}_G$ 的规范场论版本。把扩展 Bogomolny 方程的解限制到其边界值，${\EUQ}$ 便有一对到 $\EUW_0$ 的映射，换言之，一个到 $\EUW_0\times \EUW_0$ 的映射。这就是 Higgs 丛的 Hecke 对应。

${\EUQ}$ 一般而言非常复杂，具有各种不同维数的分支。然而，如果只限于 $\EUW_0$ 的一个良好区域（其中每个 Higgs 丛在模去规范变换后只有有限多个给定类型的可能 Hecke 修正），那么 ${\EUQ}$ 可视为 $\EUW_0\times \EUW_0$ 的子流形。此外，在这个良好区域中我们可以模去规范变换，${\EUQ}$ 就成为 $\MH\times\MH$ 的中间维数子流形。

't Hooft 算子保持 $(B,A,A)$ 型超对称，正如我们在第~\ref{sec:lineop}~节末所解释的。这意味着必须能把 ${\EUQ}$ 解释为 $\MH\times \MH$ 中的 $(B,A,A)$ 型膜。特别地，它必须在复结构 $I$ 中全纯，并且相对于复结构 $J$ 与 $K$——等价地，相对于式~\eqref{eq:ninto} 所定义的全纯辛形式 $\Omega_I=\omega_J+i\omega_K$——是拉格朗日的\footnote{在中间维数以上，存在不是拉格朗日的 $A$-膜 \cite{ref80}。但在良好区域中，${\EUQ}$ 是中间维数的，故必是拉格朗日的。}。更确切地说，若 $\pi_i:{\EUQ}\to\MH\times\MH$（$i=1,2$）是两个投影，则 ${\EUQ}$ 必须相对于 $\MH\times \MH$ 上的全纯辛形式 $\Omega=\pi_1^*(\Omega_I)-\pi_2^*(\Omega_I)$ 是拉格朗日的。这里的负号保证了 $\MH\times \MH$ 中的对角——即没有 't Hooft 算子的平凡 Hecke 对应——是拉格朗日的。

记 $C_+$ 与 $C_-$ 为 $W=I\times C$ 的两端，我们可以把 $\Omega$ 在 ${\EUQ}$ 上的限制写成
% 式 (10.41)
\setcounter{equation}{40}
\begin{equation}\label{eq:symplevolume}
{1\over\pi}\left(\int_{C_+}
-\int_{C_-}\right)|d^2 z|\,\Tr\left(\delta\phi_z\wedge \delta
A_\bz\right)={1\over \pi}\int_Idy\int_C|d^2z|\,
\partial_y \Tr\left(\delta\phi_z\wedge\delta A_\bz\right).
\end{equation}

为计算这个形式在 $\EUQ$ 上的限制，我们可以假定变分 $\delta\phi_z$ 与 $\delta A_\bz$ 满足 \eqref{eq:BPSextsimple}\ 的线性化。\eqref{eq:BPSextsimple}\ 中前三式的线性化给出
% 式 (10.42)
\setcounter{equation}{41}
\begin{equation}\label{eq:BPSlin}
\begin{aligned}D_\bz\delta\phi_z&=-[\delta
A_\bz,\phi_z],\\
\partial_y\delta A_\bz&=-iD_\bz\delta\phi_0 - i [\delta
A_\bz,\phi_0]\\
\partial_y\delta\phi_z&=-i[\delta\phi_z,\phi_0]-i[\phi_z,\delta\phi_0]
.\\\end{aligned}
\end{equation}
 利用这些方程、Jacobi 恒等式以及分部积分，可以验证 \eqref{eq:symplevolume}\ 为零，因此 ${\EUQ}$ 如预期的那样是复拉格朗日的。

\begin{yizhu}[原文疑误]
按：本片有一处 $z=0$ 照录原文，实应为 $z=z_0$。原文说“若 $\varphi$ 在 $z=0$ 之外给定，则越过 $z=0$ 它至多只有唯一一个全纯延拓”，但全段讨论的 't Hooft 算子位于 $z=z_0$，奇点空间也记作 $\bar {\EUY}({}^L\negthinspace w;z_0)$；紧随其后的段落写“即使 $\varphi$ 在 $y<y_0$ 时于 $z=z_0$ 处全纯……在 $y>y_0$ 时它一般也会在 $z=z_0$ 处出现极点”，可见前两处 $z=0$ 是 $z=z_0$ 的笔误。照录并标出，读者按 $z=z_0$ 理解即可，数学内容不受影响。
\end{yizhu}

\begin{yizhu}[易错点]
按：不要把 \eqref{eq:BPSstaticgen}\ 与第~\ref{sec:theck}~节的普通 Bogomolny 方程混为一谈：普通情形靠时间反演不变性可以证明 $\phi_1=A_0=0$，只剩 $W$ 上的方程；而“扩展”情形只是令 $t=1$、保留时间分量 $\phi_0$ 与 $A_0$ 作为待解的场，方程 \eqref{eq:BPSstaticgen}\ 中的 $D$、$\star$、$D^*=\star D\star$ 都是三维 $W$ 上的算子，$A_0$、$\phi_0$ 以“源”的形式出现在右端（如 $\partial_yA_{\bar z}=-iD_{\bar z}\phi_0$）。只有再借助消灭定理才可令 $\phi_1=A_0=0$，得到 \eqref{eq:BPSextsimple}。
\end{yizhu}

\begin{yizhu}[易错点]
按：本片最反直觉的一点是：丛的 Hecke 修正带参数（模空间 $\bar {\EUY}({}^L\negthinspace w;z_0)$ 的维数为 $\sum_{i<j}(m_i-m_j)$），而把 $\varphi$ 一起纳入后修正反而“消去参数”、一般只有有限多个。原因是对比的对象不同：全纯丛可以在单点被扭化而保持全纯（这正是 Hecke 修正），而全纯截面不能在单点修改还保持全纯，所以 $\varphi$ 在 $y>y_0$ 侧的全纯性对 $f_i{}^j$ 的 Taylor 系数施加了恰好 $\sum_{i<j}(m_i-m_j)$ 个条件——数目与丛的修正空间维数相同。当 $\varphi(z_0)$ 正则半单时该向量场零点有限，故可能的 Hecke 修正构成有限集。
\end{yizhu}

\begin{yizhu}[物理图像]
按：可把 $\varphi$ 想成 Schubert 簇 $\bar {\EUY}({}^L\negthinspace w,z_0)$ 上的一个向量场，其零点就是“允许的”修正；这比逐条数 Taylor 系数更几何。更关键的物理后果是特征多项式与 $y$ 无关，即使越过 't Hooft 算子也不变，故 $\varphi$ 所在的 Hitchin 纤维被保持：缠在 Hitchin 纤维上的膜（带平坦 Chan-Paton 丛者）在 Hecke 变换下支撑不变。这正是“磁特征膜”的直观内容，也是本片末尾要在 $\MH\times\MH$ 上构造 Hecke 对应 ${\EUQ}$ 的动机。
\end{yizhu}

\begin{yizhu}[易错点]
按：$\Omega=\pi_1^*(\Omega_I)-\pi_2^*(\Omega_I)$ 中的负号不是任意约定。限制到对角 $\Delta\subset\MH\times\MH$ 上，两项相消给出 $0$，于是平凡 Hecke 对应（无 't Hooft 算子）的支撑是拉格朗日的；若换成加号，拉回将是 $2\Omega_I$，对角就辛而非拉格朗日。另外 $(B,A,A)$ 型的读法是：在复结构 $I$ 中全纯（$B$ 的首字母），在 $J$、$K$ 中均为 $A$-型（即相对于 $\Omega_I=\omega_J+i\omega_K$ 拉格朗日）；\eqref{eq:symplevolume}\ 的验证只用到 \eqref{eq:BPSextsimple}\ 的线性化、Jacobi 恒等式与分部积分，故 ${\EUQ}$ 在该良好区域确实是复拉格朗日的。
\end{yizhu}
